\begin{recipe}{Shoarmabroodjes}
\kicker{Brood \& basis · zelf bakken}
\index[register]{Shoarmabroodjes}
\index[register]{Bloem}

\heroplaceholder{handgetekende illustratie van opgebolde broodjes op de steen}
\meta{2,5 uur rijzen \dvd 7 stuks \dvd Midden-Oosters}

\lettrine{Z}{achte,} lichtjes gezwollen broodjes om zelf te vullen met shoarma.
Het deeg is eenvoudig, het enige wat het vraagt is tijd voor het rijzen.
Een hete pizzasteen zorgt voor de goede bite aan de onderkant.

\blockrule{Ingrediënten}
\begin{ingredients}
  \ing{300 g}{bloem}
  \ing{360 ml}{lauwwarm water}
  \ing{32 g}{suiker}
  \ing{8 g}{zout}
  \ing{5 g}{instant gist}
\end{ingredients}

\blockrule{Bereiding}
\tip{Hoe langer het deeg rijst, hoe meer smaak. Twee uur op kamertemperatuur werkt goed; langer mag ook.}
\ingredientsketch{broodjesdeeg}
\begin{steps}
  \item Meng bloem, suiker, zout en gist in een grote kom. Voeg het lauwe water
        toe en meng tot een deeg.
  \item Kneed 5–8 minuten door tot het deeg glad en elastisch aanvoelt.
  \item Dek de kom af en laat het deeg 2–3 uur rijzen op kamertemperatuur.
  \item Zet een pizzasteen in de oven en verhit minstens 30 minuten op 250\,°C.
  \item Sla de lucht uit het deeg en verdeel in porties van 100\,g. Bol elke
        portie op en laat 15 minuten afgedekt rusten.
  \item Rol elk bolletje uit tot een schijf van ca.\ 16\,cm.
  \item Leg de broodjes op bakpapier en schuif ze op de hete steen. Ze zijn
        klaar zodra ze opbollen, na een paar minuten.
\end{steps}
\end{recipe}
